\subsection{Deformations due to heat load}


Synchrotron radiation sources produce extremely intense beams of photons across a broad spectrum, and managing the \textbf{heat load} deposited by this radiation is one of the central engineering challenges in beamline design.

Heat load refers to the power deposited by the synchrotron radiation beam onto optical components such as masks, slits, absorbers, monochromators, and mirrors. Modern storage rings and especially fourth-generation sources (such as EBS at ESRF, APS-U, or SPring-8-II) produce beams with total powers ranging from \textbf{hundreds of watts to tens of kilowatts}, concentrated into very small footprints.

The critical parameter is not just total power but also the peak power density: the power per unit area that strikes an optical surface.


\paragraph{Overview of the Simulation Loop}

The simulation pipeline follows this chain:

\begin{itemize}
    \item Source power spectrum calculation OASYS/XOPPY
    \item Profile of incident power on the optical element (taken into account the effect of the elements upstream of it)
    \item Temperature field (FEA)
    \item Surface deformation (FEA)
    \item Extract main parameters of deformation (e.q. maximum slope error)
    \item Iteration and optimization of the main parameters (cooling, dimensions, shapes) using FEA
    \item Extract FEA-calculated surface for ray tracing verification of the optical effect of the residual deformations
    \item Optical output (flux, focus, resolution)
\end{itemize}

In this chapter we study for a monochromator crystal the broadening of the energy distibution due to the surface deformation and the image spot profile degraded by thermal deformation.
